% Lösungen Flächeninhalt durch Integration (zusammenbau v0.4, Heft)
\einheitenkopf{Flächeninhalt durch Integration}

\erg{8}{a) Integral – der obere Rand ist gekrümmt \quad b) Integral $\frac{8}{3}$ plus Dreieck $4$; $A = \frac{20}{3} \approx 6{,}67$ \quad c) $3 + \frac{8}{3}$; $A = \frac{17}{3} \approx 5{,}67$}
\erg{9}{a) (1) $y = 2x + 2$ ist die Sekante durch die Nullstelle $-1$ und den y-Achsenabschnitt; (2) Segment zwischen Graph und Sekante $= \frac{1}{6}$; (3) Dreieck unter der Sekante mit den Achsen; (4) das Segment ist ein Sechstel des Dreiecks. \quad b) (1) $y = 3x + 3$ ist die Sekante durch die Nullstelle $-1$ und den y-Achsenabschnitt; (2) Segment zwischen Graph und Sekante $= \frac{1}{6}$; (3) Dreieck unter der Sekante mit den Achsen $= \frac{3}{2}$; (4) das Segment ist ein Neuntel des Dreiecks. \quad c) Tangente $t(x) = -2x + 5$, Nullstelle $2{,}5$; Dreieck unter $t$: $\frac{1}{2} \cdot 1{,}5 \cdot 3 = 2{,}25$; unter $f$ von $1$ bis $2$: $\frac{5}{3}$; Querschnitt $= \frac{7}{12}\,\text{m}^2$; Volumen $= \frac{7}{12} \cdot 6 = 3{,}5\,\text{m}^3$ \quad d) Tangente $t(x) = -x + 2{,}5$; Dreieck unter $t$: $\frac{1}{2} \cdot 1{,}5 \cdot 1{,}5 = 1{,}125$; unter $f$ von $1$ bis $2$: $\frac{5}{6}$; Querschnitt $= \frac{7}{24}\,\text{m}^2$; Volumen $= \frac{7}{3} \approx 2{,}33\,\text{m}^3$}
\erg{10}{a) vorwärts (Inhalt gesucht) – Grenzen bestimmen, integrieren \quad b) Schnittstellen $0$ und $\frac{m}{2}$; Inhalt $\frac{m^3}{24} = 9$; $m = 6$ \quad c) Gesamt $8$; $k^3 = 4$; $k = \sqrt[3]{4} \approx 1{,}59$}
\erg{11}{a) $k = -1$: Das Intervall von $-1$ bis $1$ liegt symmetrisch zu $0$; das Rechteck mit Höhe $3$ und Breite $2$ hat den Inhalt $6$, die Flächenstücke über und unter der Geraden $y = 3$ sind wegen der Punktsymmetrie gleich groß und heben sich auf. \quad b) $k = -2$: Das Intervall von $-2$ bis $2$ liegt symmetrisch zu $0$; das Rechteck mit Höhe $5$ und Breite $4$ hat den Inhalt $20$, die Flächenstücke über und unter der Geraden $y = 5$ sind wegen der Punktsymmetrie gleich groß und heben sich auf.}
\erg{12}{a) Wert negativ – nicht die Fläche, die Fläche ist sein Betrag \quad b) Wert $= 8$ (Trapez unter dem Graphen, ganz oberhalb der Achse) \quad c) Das Stück links liegt über, das Stück rechts unter der x-Achse; beide sind wegen der Punktsymmetrie zum Ursprung gleich groß, ihre orientierten Inhalte heben sich auf.}
\erg{13}{a) Die Stücke zwischen Graph und Gerade links und rechts von $W$ sind wegen der Punktsymmetrie gleich groß und liegen auf verschiedenen Seiten der Geraden; das Integral ist also die Fläche unter $g$, das Dreieck mit Grundseite $4$ und Höhe $g(5) = f(5)$: Wert $= \frac{1}{2} \cdot 4 \cdot 4 = 8$ \quad b) Die Stücke zwischen Graph und Gerade links und rechts von $W$ sind wegen der Punktsymmetrie gleich groß und liegen auf verschiedenen Seiten der Geraden; das Integral ist also die Fläche unter $g$, das Dreieck mit Grundseite $6$ und Höhe $g(7) = f(7)$: Wert $= \frac{1}{2} \cdot 6 \cdot 6 = 18$}
